\begin{corollary}[Proved under stated hypotheses: mixed mild surface boundary]
\label{D:Gorenstein-lc-boundary}
Let $S$ be a normal proper connected surface over an algebraically
closed field, and let $Y\subset S$ be finite.  Suppose that every
point of $Y$ is regular or belongs to at least one of the following
classes:
\begin{enumerate}
\item a log terminal surface singularity;
\item a singularity admitting a resolution whose reduced exceptional
fibre is one smooth elliptic curve, with no $\mathbb Q$-Gorenstein
assumption;
\item a Gorenstein log canonical cusp singularity, whose minimal
exceptional fibre is a rational nodal cycle.
\end{enumerate}
Then $\pi_1(S\setminus Y)$ is topologically finitely generated.
In particular, this holds if every omitted singularity is either
log terminal or Gorenstein log canonical.  The singularities
retained in $S\setminus Y$ are unrestricted.
For a possibly nonnormal proper scheme of dimension at most two,
the same conclusion holds if the finite inverse image of $Y$ on
each normal surface component satisfies these conditions and $Y$
has codimension at least two in every irreducible component.
\end{corollary}
\begin{proof}
Choose a regular resolution $q:X\to S$ which is minimal near the
omitted log terminal and cusp points, is the specified single-curve
resolution near the omitted elliptic points, and is an isomorphism
at the omitted regular points.  At an irreducible nodal cusp core,
blow up its node once.  Resolve the singularities retained in the
open arbitrarily.  As in Corollary~\ref{D:elliptic-surface-boundary},
these finitely many local choices glue along the regular locus.
The regular proper surface $X$ is projective by
\cite[Lemma 54.16.11, tag 0C5P]{StacksDSurfaceResolution}.
Put $U=S\setminus Y$ and $U'=q^{-1}(U)$.

Every connected finite \'etale cover of $U$ stays connected after
pullback to $U'$: its generic fibre is a field and the pullback is
normal, with every connected component dominating the integral
base $U'$.  Its normalization over $S$ is finite and quasi-\'etale.
At every divisorial boundary component over an elliptic point,
Proposition~\ref{D:single-curve-residue-tameness} proves tameness.
At a cusp point, Theorem~\ref{D:cusp-tame} proves tameness, including
at the components introduced by the node blowup.  Thus the image
of wild inertia in $\pi_1(U)$ is zero at all these components.

At an omitted log terminal point, the punctured strict-local
fundamental group is finite by
\cite[Theorem 5.3, p.~739; see also Corollary 6.4, p.~741]
{CarvajalRojasYasuda2025B}.  The theorem holds in every
characteristic, including two and three.  The wild inertia image
at an exceptional divisor over this point factors through the
image of that local group in $\pi_1(U)$, and is therefore finite.
Regular omitted points give only codimension-two boundary on
$X$ and require no divisorial condition.
Consequently all boundary wild images on this fixed projective
modification are finite.  There is no original divisorial
boundary, so $Q(S,Y)=\pi_1(U)$, and
Proposition~\ref{D:finite-exceptional-wild} proves the conclusion.

Lemma~\ref{D:Gorenstein-lc-configurations} proves the asserted
Gorenstein log canonical special case: its all-zero configuration
is canonical, hence log terminal, and its other configurations
are the elliptic and cusp cases just treated.
In the nonnormal case we pass to the finite disjoint union of
normalizations of the reduced irreducible components.  The proved
normal case applies to each surface component, while components
of dimension at most one have no omitted point and are proper.
Finite proper descent, as in
Proposition~\ref{D:fixed-boundary-positive}, gives the conclusion
downstairs.
\end{proof}
